\section{Conclusion}

This study developed a paired wavelet attenuation--coherence framework to evaluate timescale-specific PM$_{2.5}$ buffering by urban parks in Bangkok using defensible 2021 hourly park--outside station pairs. The central finding is that urban parks did not act as uniform mean-concentration reducers; rather, they modified PM$_{2.5}$ variability in distinct ways across timescales. Attenuation was negative overall, with the strongest damping generally occurring from sub-diurnal to multi-day bands, whereas coherence increased toward the diurnal and 2--7 d bands, indicating that parks often reduced fluctuation amplitude while remaining embedded within broader city-scale temporal rhythms.

A second key result is that Bangkok parks were not interchangeable. The retained pairs organized into three interpretable regimes---strong decouplers, moderate broadband dampers, and weak followers or mixed cases---showing that park air-quality function depends not only on whether variability is reduced, but also on how strongly park dynamics remain coupled to surrounding urban conditions. This makes the attenuation--coherence framework more informative than static inside--outside mean comparison alone.

The explanatory analysis further indicated that average greenness was not sufficient to explain this heterogeneity. Mean land-cover contrast provided a useful first-order signal, annulus-based land-cover gradients provided the broadest explanatory improvement, water-related context contributed strong but more targeted coherence-oriented associations, and reviewed morphology remained secondary. Together, these findings suggest that PM$_{2.5}$ buffering depends more on layered spatial structure---including protected vegetated interiors, surrounding built-up pressure, comparator-site context, and supportive local microenvironment---than on total green area alone.

From a practical perspective, the study suggests that urban park design should move beyond area-based greening targets and give greater attention to protected interior vegetation, reduced central hardscape pressure, stronger park-edge transitions, and the immediate surrounding urban perimeter. More broadly, the study provides a reproducible framework for evaluating how urban green spaces interact with air-pollution dynamics in time, not only in mean concentration space. The next priority is replication with denser multi-year monitoring and stronger directional context so that the Bangkok regime structure can be tested, refined, and translated more confidently into urban environmental design.